\chapter{Conclusion}
\label{sec:conclusion}

\section{Overview}

This thesis studied the design of Bayesian two arm trials with a normal endpoint that borrow
historical information on the control arm. Borrowing makes a trial smaller and reduces the number of
patients, but it costs the frequentist guarantee that the type one error (T1E)
rate is controlled at every value of the control parameter. 

\citet{best_beyond_2025} restore that guarantee by integrating the conditional T1E
over a design prior for the control parameter to obtain the \textit{average T1E rate} (avgT1E). An
analogous \textit{average power} (avgPower) was developed here, and closed form expressions and
properties were derived for both under normal priors, together with the sample size computation they
imply. Because a single normal analysis prior offers no protection when the historical data disagree
with the new trial, the analysis prior was then built as a \textit{robust meta-analytic-predictive}
(MAP) prior \citep{schmidli_robust_2014}, for which the same metrics follow by numerical integration.

On top of these two metrics the \textit{Experimental Unit Information Index} (EUII)
\citep{held2025balancing} was built in average form. It normalizes the \textit{diagnostic odds ratio}
(DOR) by the number of patients a design needs, placing designs of different size, and designs whose
size is random, on a common per patient scale. The framework was extended to \textit{dual criterion}
(DC) designs, which additionally require the effect to exceed a clinically relevant threshold, and to
Bayesian group sequential designs, which stop early for efficacy or for futility. Since adding interim
looks inflates the avgT1E, a calibration algorithm was developed to place every design at a common
avgT1E and avgPower before its EUII is read. A Phase II Crohn's disease trial served as the running
example throughout the methodology and as the case study.


\section{Methodological Contributions}

\paragraph{An average power}
Following from the avgT1E, an analogous avgPower, obtained by integrating
the conditional power over the design prior, was defined
(Section~\ref{subsec:metrics_bayesian}). Together, avgT1E and avgPower were used to build the EUII in
average form.

\paragraph{Closed forms and three properties}
Closed form expressions for avgT1E and avgPower were derived under normal priors
(Section~\ref{subsec:closedforms_normal}), and three
properties were established. Average metrics are always pulled closer to $0.5$ than their conditional counterparts.
Under aligned analysis and design priors, avgT1E is controlled exactly at the nominal level. And as
the design prior becomes vague, both avgT1E and avgPower converge to $0.5$, whatever the analysis
prior.

\paragraph{Sample size computation}
A closed form for the sample size needed to reach a target average power was obtained for a single
normal analysis prior (Section~\ref{subsec:samplesize_closedform}). For the MAP prior and its robust
extension no closed form exists, but since
avgPower is monotone in $n_C$, numerical inversion gives the smallest adequate sample size.

\paragraph{Normalization of the index}
Under borrowing, it is not obvious whether the patients encoded in the analysis prior should enter the
sample size that normalizes the EUII alongside those enrolled in the trial
(Section~\ref{subsec:euii}). We argued they should not:
the index answers how much evidence each newly enrolled patient contributes, historical patients are
already collected, and including them would conflate the choice of prior with the choice of trial
size, penalizing an informative prior for the very information it supplies.

\paragraph{Asymptotic behavior}
The asymptotic value of the EUII was derived under the single normal model for both criteria
(Section~\ref{subsubsec:asymptotic_value}), giving a
finite limit as the effect size grows and a reference for reading finite sample behavior.

\paragraph{The framework does not extend to contrast borrowing}
The framework does not transfer to borrowing on the treatment contrast, for three reasons
(Section~\ref{subsec:contrast_limitation}). A prior on
$\delta$ typically favors a beneficial effect and so conflicts with the null, forcing strict
frequentist error control to discard the historical information. No exact analogue of the avgT1E
control result exists in that setting. And the EUII itself cannot be built, since the
metric proposed for this case, a Type III error, is a joint probability, not a conditional
one.

\section{Findings from the Application}

\paragraph{The robustness weight increases tolerance to drift, and location matters more than spread}
Varying the design prior by switching, shifting its location, and scaling its spread
(Section~\ref{sec:fixed_design}) showed that only
a control mean better than $\pi_{\text{MAP}}$ assumes can inflate avgT1E, and that being wrong about
where $\theta_C$ lies costs considerably more than being wrong about how uncertain it is. Raising
the robustness weight reduces avgT1E under every form of conflict considered.

\paragraph{The EUII is not a diagnostic for prior misspecification}
Under misaligned priors, EUII stays close to its aligned value
even though avgT1E and avgPower move substantially, because the index is a ratio of odds and a
disagreeing prior shifts both in the same direction, so the shifts largely cancel
(Section~\ref{sec:fixed_design}). A badly optimistic
design and a badly underpowered one can therefore achieve nearly the same EUII. Prior sensitivity
should be assessed through avgT1E and avgPower, not EUII, which should be used for comparing
designs under a correctly specified prior.

\paragraph{EUII for DC versus EUII for SC}
The EUII curves of the two criteria cross: the DC has the higher EUII for effects moderately above
the decision value, the SC for large effects. Raising the success threshold lowers avgT1E by a
fixed amount but grows the loss in avgPower with $\delta$, so the fixed gain dominates at moderate
effects and the loss dominates at large ones, a crossing that also appears as a function of the
control sample size (Section~\ref{sec:fixed_design}). 

\paragraph{Interim analyses and futility raise the evidentiary value per patient}
Every sequential design has a higher EUII than its fixed counterpart at the same $\delta$, and a second
interim raises it further (Sections~\ref{sec:one_interim} and~\ref{sec:two_interim}) but inflates the avgT1E. 
Adding a futility rule raises the EUII while leaving avgT1E and avgPower almost untouched, since it mostly remove
trials that would have failed at the final analysis anyway.


\paragraph{A calibration algorithm for Bayesian group sequential designs}
Comparing sequential designs at a fixed nominal threshold is not always a fair comparison, since designs 
with different numbers of looks or different stopping rules would then have different avgT1E 
and avgPower. A calibration algorithm was therefore developed that, for a chosen stopping rule 
and number of looks, tunes the threshold together with the maximum sample size to reach a 
specified avgT1E and avgPower (Section~\ref{sec:seq_calibration}).
At a fixed target, more interim looks always gave a higher EUII (Section~\ref{sec:calibrated_sc}), and 
among stopping rules at a fixed number of looks, Pocock's style constant threshold almost always gave the highest EUII (Section~\ref{sec:calibrated_shapes}).



\section{A Design Workflow}
\label{sec:workflow}
The pieces developed in this thesis compose into a fixed order for planning a new trial in this framework.
\begin{enumerate}
    \item \textbf{Build $\pi_{\text{MAP}}$} from the historical arms, including only those studies that satisfy prespecified
    eligibility criteria, for example the ones defined by \citet{pocock_combination_1976}.
    \item \textbf{Choose $\omega$}, fixing the analysis prior $\pi_A$ and, with it, $n_A = \mathrm{ESS}(\pi_A)$.
    \item \textbf{Set the threshold} $p = 1 - \alpha^\ast$. No computation is required: exact under alignment, and, as shown in Section~\ref{sec:fixed_design}, conservative under $\pi_D = \pi_{\text{MAP}}$ for every $\omega > 0$.
    \item \textbf{Solve for $n_C$} at the target power under $\pi_D = \pi_{\text{MAP}}$, then $n_T = r_T n_C$.
    \item \textbf{Stress test} $\omega$ against a declared \textit{drift}, by shifting $\pi_{\text{MAP}}$. If the worst case fails the desired tolerance, raise $\omega$ and return to step ii.
\end{enumerate}
A mixture analysis prior changes step iv: control of avgT1E still holds under alignment,  but the closed 
form of Section~\ref{subsec:samplesize_closedform} is no
longer available and $n_C$ follows from a bisection on avgPower, which is monotone in $n_C$
(Section~\ref{sec:fixed_design}).

A sequential design changes step iii as well. Adding interim looks inflates the avgT1E
and the threshold $p$ must be calibrated by root finding (Algorithm~\ref{alg:calibrate_threshold}). 
Since that threshold depends on the schedule and the
schedule on it in turn, steps iii and iv stop being separate and merge into the nested searches of
Algorithm~\ref{alg:calibrate_design}, which recalibrates the threshold at every candidate size
(Section~\ref{sec:seq_calibration}).

\section{Future Work}


 A direction for future work is extending the framework to borrowing on the treatment contrast $\delta$
directly, rather than on the control arm, addressing the issues that
Section~\ref{subsec:contrast_limitation} leaves open. 


A further direction is extending the endpoint beyond the normal case with known standard
deviation $\sigma$ used throughout the thesis, or investigating the consequences of misspecifying
$\sigma$. Moreover, the only priors considered are the MAP prior and its robust extension, but it
would be interesting to test the framework using power priors and commensurate priors.



The sensitivity analyses obtained by changing the design prior could also extend past the three
perturbations reported here, for instance by reweighting the mixture components of
$\pi_{\text{MAP}}$ or refitting it with one historical study left out at a time.

Regarding Bayesian group sequential design, a more accurate comparison of EUII using calibrated DC
designs, or varying the priors, could be performed, together with a futility boundary following a stopping rule of
its own, calibrated jointly with the efficacy boundary. The timing of the interim analyses could
likewise be optimized against the index rather than inherited from the protocol.

Finally, Bayesian designs with robust borrowing remain underused in practice, but their uptake is
expected to increase following the recent FDA guidance ``Use of Bayesian Methodology in Clinical
Trials of Drug and Biological Products: Draft Guidance for Industry'' \citep{fda_bayesian_2026}. As
these designs move into wider use, it becomes more important to understand how much evidentiary value they
deliver, how sensitive avgT1E and avgPower are to the design prior they are evaluated under, and how much prior-data conflict the
robustness weight $\omega$ can tolerate. The conclusions drawn here rest on a single motivating trial, so applying the design
workflow of Section~\ref{sec:workflow} prospectively to new trials, or retrospectively to completed
ones, would show how far they generalize.


%%% Local Variables:
%%% mode: latex
%%% TeX-master: "../MasterThesisSfS"
%%% End:
